\documentclass[12pt, a4paper]{academics}

\academicsCoverLabel{解题报告}
\academicsCourse{算法设计与分析}{Design and Analysis of Algorithms}
\academicsReport{网络流}{2026 年 6 月 7 日}
\academicsArthur{华博文}{19240212}

\begin{document}

\makeacademicscover

\section{引言}

网络流是图论与组合优化中最为深邃的领域之一。其核心问题——在容量受限的有向网络中如何最大化从源点到汇点的流量——表面上是一个约束优化问题，但在 Ford 和 Fulkerson 于 20 世纪 50 年代提出最大流最小割定理之后，它成为了连接组合数学、线性规划与实际工程问题的桥梁。从交通调度、通信路由到图像分割、棒球淘汰问题，最大流及其变体在无数实际场景中起着基础性的作用。

本报告以网络流为主题，选取四道难度递进、思想关联的题目，系统展示从基础最大流算法到最小费用流扩展的完整技术脉络。题一以 Edmonds--Karp 算法求解基础最大流问题，引入残量网络与 BFS 增广的核心思想；题二将数据规模提升两个数量级，迫使我们在层次图、多路增广与当前弧优化的配合下实现 Dinic 算法，将复杂度从 $\mathrm{O}(VE^2)$ 降至实用范围；题三将最大流转化为组合工具，对二分图最大匹配问题建立流网络模型，并借助 K\"onig 定理从残量网络中提取最小点覆盖；题四则引入费用维度，以势能法配合 Dijkstra 实现 SSP（Successive Shortest Path）算法，在保证最短路计算效率的同时处理残量网络中不可避免的负权边。

四道题目犹如网络流这一宏大主题的四块积木：从增广路到层次图，从单纯最大流到组合应用，再到费用与流量的联合优化，每一步都在前一步的基础上引入新的概念与技巧，而每一步也都反过来加深了对前一步的理解。本报告中的所有代码均使用 C++23 标准编写，在 Apple clang 21.0.0 编译器下通过测试。

\section{题一：Edmonds--Karp 算法与基础最大流}

\subsection{问题重述}

给定一个有向网络，包含 $n$ 个顶点和 $m$ 条有向边，每条边 $(u, v)$ 上带有一个非负容量 $c(u, v)$。指定源点 $s$ 和汇点 $t$，要求从 $s$ 到 $t$ 输送尽可能多的流量，满足以下约束：对于任意顶点 $v \notin \{s, t\}$，流入量与流出量相等（流量守恒）；对于任意边 $(u, v)$，通过该边的流量不超过其容量；所有边的流量均为非负值。最终输出从 $s$ 到 $t$ 的最大流量的值。

输入格式为：第一行四个整数 $n, m, s, t$，分别表示顶点数、边数、源点和汇点（$1 \le s, t \le n$）。接下来 $m$ 行每行三个整数 $u_i, v_i, c_i$，表示一条从 $u_i$ 到 $v_i$ 的边，容量为 $c_i$。数据范围：$2 \le n \le 500$，$1 \le m \le 1000$，容量不超过 $10^9$。


\subsection{残量网络与最短增广路策略}

最大流问题的经典求解框架是 Ford--Fulkerson 方法：从一个零流开始，反复在残量网络中寻找一条从源点到汇点的增广路，沿该路径增加流量，直至不存在增广路为止。残量网络以原始网络为基础，对于每条边 $(u, v)$，残量网络中同时保留正向剩余容量 $c(u, v) - f(u, v)$ 和反向容量 $f(u, v)$，后者对应于「撤销」已输送流量的能力。这一设计的核心洞察在于：每当我们沿某条边输送了流量，我们立即在反向创建一条等量的容量，使得后续的增广操作可以「后悔」此前的决策，重新调配流量分布。

Ford--Fulkerson 方法并不指定如何选择增广路——若随意选择，在边容量为无理数时甚至可能不终止。Edmonds 和 Karp 的关键贡献在于规定每次选取边数最少的增广路，即使用广度优先搜索（BFS）在残量网络中寻找最短路径。这一策略保证了每次增广路的长度单调不降，从而将增广次数的上界限制为 $\mathrm{O}(VE)$（每条边至多成为瓶颈 $V$ 次），每次 BFS 耗时 $\mathrm{O}(E)$。实现时，BFS 从源点出发，沿残量网络中容量为正的边层层推进，直至抵达汇点，沿途记录每个顶点的父顶点及所经过的边的索引，以便回溯瓶颈容量和更新残量网络。

Edmonds--Karp 算法的时间复杂度为 $\mathrm{O}(VE^2)$。具体而言，每次 BFS 扫描所有边，耗时 $\mathrm{O}(E)$；增广路径的长度单调递增且不超过 $V-1$，每条边成为关键边（路径中容量最小的边）的次数不超过 $V/2$，因此总增广次数为 $\mathrm{O}(VE)$，二者相乘即得 $\mathrm{O}(VE^2)$。对于 $n \le 500$、$m \le 1000$ 的数据规模，最坏情况下约 $2.5 \times 10^8$ 次操作，完全可以在时限内完成。空间复杂度为 $\mathrm{O}(V + E)$，用于存储邻接表和 BFS 辅助数组。

\subsection{代码实现}

\inputminted{c++}{../src/p1_edmonds_karp.cpp}

\section{题二：Dinic 算法与大规模最大流}

\subsection{问题重述}

本题的问题描述与题一完全一致——给定有向网络，求从源点 $s$ 到汇点 $t$ 的最大流。唯一的变化在于数据范围急剧扩大：$2 \le n \le 10000$，$1 \le m \le 100000$。若仍沿用 Edmonds--Karp 的 $\mathrm{O}(VE^2)$ 算法，在最坏情况下将达到 $10^{13}$ 次操作，完全不可接受。这意味着我们必须从算法结构层面进行根本性的改进，而非仅仅在常数上做文章。


\subsection{层次图与阻塞流的批量推送}

Edmonds--Karp 的核心瓶颈在于「一次增广只用一条路径」——每次 BFS 的成果仅被利用一次便告废弃，下一轮增广又需重新 BFS。Dinic 算法的核心洞见是：将一次 BFS 所揭示的层次结构复用多次，在一次 BFS 构建的层次图中「榨干」所有可用的增广流量，此即所谓的阻塞流（blocking flow）。这样一来，BFS 的调用次数从 $\mathrm{O}(VE)$ 骤降至 $\mathrm{O}(V)$（因为每发送一轮阻塞流，源点到汇点的最短距离严格增加），总体复杂度降至 $\mathrm{O}(V^2E)$。

Dinic 算法由两个交替执行的阶段构成。第一阶段是 BFS 构建层次图（level graph）：从源点出发，仅沿残量网络中容量为正的边前进，为每个顶点分配一个层次 $\mathit{level}[v]$（表示从源点到 $v$ 的最短距离）。若汇点不可达，则算法终止。第二阶段是 DFS 发送阻塞流：从源点开始，沿着满足 $\mathit{level}[v] = \mathit{level}[u] + 1$ 的边进行深度优先搜索，不断寻找通往汇点的路径并增广，直至层次图中不再存在任何从源点到汇点的路径。这种「一轮 BFS、多轮 DFS」的批处理模式是 Dinic 相较 Edmonds--Karp 的本质飞跃——它将「逐条增广」的线性思维切换为「分层榨取」的批量思维。

\subsection{当前弧优化与多路增广}

单纯的 DFS 回溯式增广仍然可能反复访问同一条已耗尽的边。Dinic 引入了当前弧优化：为每个顶点维护一个指针 $\mathit{ptr}[u]$，记录该顶点尚未被完全利用的第一条出边的索引。在 DFS 过程中，每当我们从某条边出发无法推送任何流量时，就将指针后移，永久跳过这条「死边」。这一优化至关重要——它将每条边在每一轮阻塞流中的访问次数限制为常数次，从而将单轮 DFS 的复杂度控制在 $\mathrm{O}(VE)$（最坏情况），实际上在多路增广的配合下效率远高于此。

多路增广则意味着 DFS 不会在找到单条增广路后立即返回，而是在递归返回途中继续尝试其他分支（代码中通过 \mil{return tr} 保持流量累加）。这进一步减少了 DFS 的重复调用次数。配合当前弧优化后，每条边在每轮阻塞流中至多被考虑一次，DFS 的整体开销被牢牢控制在合理范围内。

Dinic 算法的一般时间复杂度为 $\mathrm{O}(V^2E)$。每轮 BFS 构建层次图耗时 $\mathrm{O}(E)$，BFS 至多被调用 $\mathrm{O}(V)$ 次（因为汇点的层次严格递增），每轮 DFS 发送阻塞流的复杂度为 $\mathrm{O}(VE)$——每条边在每轮 DFS 中至多被考虑一次（当前弧优化的保证）。对于 $n = 10000$、$m = 100000$ 的数据规模，理论上界为 $10^{13}$ 量级，但在实际随机网络中，层次递增的速度远快于理论最坏情形，实际运行效率极高。在二分图匹配等特殊场景中（所有边容量为 $1$），Dinic 的复杂度可进一步降至 $\mathrm{O}(E\sqrt{V})$。空间复杂度为 $\mathrm{O}(V + E)$。

\subsection{代码实现}

\inputminted{c++}{../src/p2_dinic.cpp}

\section{题三：二分图最大匹配与最小点覆盖}

\subsection{问题重述}

给定一个二分图，左侧有 $n$ 个顶点，右侧有 $m$ 个顶点，图中共有 $k$ 条边（每条边连接左侧某顶点与右侧某顶点）。要求计算该二分图的最大匹配的大小，即最多能够选出多少条边，使得其中任意两条边都不共享端点。此外，还需输出一个最小点覆盖——即一个顶点集合，使得图中的每条边至少有一个端点属于该集合，且该集合的大小最小。由 K\"onig 定理，二分图的最大匹配大小等于最小点覆盖的大小，这为从最大匹配推导最小点覆盖提供了理论依据。

输入格式：第一行三个整数 $n, m, k$（$1 \le n \le 1000$，$1 \le m \le 1000$，$1 \le k \le 10000$）。接下来 $k$ 行每行两个整数 $u_i, v_i$，表示左侧顶点 $u_i$ 与右侧顶点 $v_i$ 之间有一条边（$1 \le u_i \le n$，$1 \le v_i \le m$）。输出第一行为最大匹配的大小，第二行为最小点覆盖的大小，第三行为最小点覆盖中的顶点列表（左侧顶点用 \mil{L} 前缀，右侧顶点用 \mil{R} 前缀）。


\subsection{流网络建模：二分图到最大流的归约}

二分图最大匹配可以直接建模为最大流问题，这一转化是网络流模型威力的一次生动展示。构造一个流网络如下：新增一个超级源点 $S$ 和一个超级汇点 $T$。从 $S$ 向每个左侧顶点连一条容量为 $1$ 的边（确保每个左侧顶点至多被匹配一次）；从每个右侧顶点向 $T$ 连一条容量为 $1$ 的边（确保每个右侧顶点至多被匹配一次）；对于原图中的每条边 $(u, v)$（$u$ 在左侧，$v$ 在右侧），连一条从 $u$ 到 $v$ 的容量为 $1$ 的边。在该网络上运行最大流算法，流量的值即为最大匹配的大小，饱和的「左-右」边即对应匹配边。

这一归约之所以成立，其根本原因在于：最大流的流量守恒约束在二分图结构中恰好对应匹配的「一端至多一匹配」约束。超级源汇上的单位容量边限制了每个左侧顶点至多发出 $1$ 单位流量、每个右侧顶点至多接收 $1$ 单位流量，而中间边上的容量 $1$ 与流量整数性保证了每条边至多被选中一次。于是，整数最大流与最大匹配之间建立了一一对应关系。

由于该网络中所有边的容量均为 $1$，Dinic 算法在此类单位容量网络上的时间复杂度退化为 $\mathrm{O}(E\sqrt{V})$，远优于一般图的 $\mathrm{O}(V^2E)$。这是因为每轮阻塞流至少会使层次增加并且在单位容量下层次上限为 $\mathrm{O}(\sqrt{V})$——直觉上，经过 $\sqrt{V}$ 轮层次图之后，剩余的增广路长度必然超过 $\sqrt{V}$，而此时剩余可增广流量已经十分有限。建模后的流网络包含 $n + m + 2$ 个顶点和 $n + m + k$ 条边，容量均为 $1$。在 $n, m \le 1000$、$k \le 10000$ 的条件下，顶点数与边数均在 $10^4$ 量级，算法在数十毫秒内即可完成。空间复杂度为 $\mathrm{O}(V + E)$。

\subsection{K\"onig 定理与最小点覆盖的构造性提取}

最小点覆盖的提取利用了 K\"onig 定理的构造性证明。在最大流结束后，考察残量网络：从源点 $S$ 出发，沿残量网络中容量为正的边进行 BFS，标记所有可达的顶点。定义集合 $X$ 为所有未被标记的左侧顶点，集合 $Y$ 为所有被标记的右侧顶点。则 $X \cup Y$ 构成一个最小点覆盖，其大小恰好等于最大匹配的大小。

这一构造的直观含义是：$X$ 中的左侧顶点对应源点 $S$ 到它们的边已被流完全占用、且无法在残量网络中被回溯的顶点；$Y$ 中的右侧顶点则是在残量网络中仍然可以从未匹配的左侧顶点或通过反向边到达的顶点，它们必须被选中以「阻挡」交替路径进一步延伸。整组操作的额外开销仅为一次 $\mathrm{O}(E)$ 的 BFS，不增加渐进复杂度。

\subsection{代码实现}

\inputminted{c++}{../src/p3_bipartite.cpp}

\section{题四：最小费用最大流}

\subsection{问题重述}

给定一个有向网络，每条边 $(u, v)$ 上除了容量 $c(u, v)$ 之外，还带有一个单位流量的费用 $w(u, v)$。指定源点 $s$、汇点 $t$ 和需求流量 $f$，要求从 $s$ 到 $t$ 输送恰好 $f$ 单位的流量，使得总费用最小——总费用定义为每条边上的流量乘以其单位费用的和。如果无法输送恰好 $f$ 单位的流量（即最大流 $< f$），则输出 $-1$。

输入格式：第一行五个整数 $n, m, s, t, f$。接下来 $m$ 行每行四个整数 $u_i, v_i, c_i, w_i$，分别表示边的起点、终点、容量和单位费用。数据范围：$2 \le n \le 500$，$1 \le m \le 5000$，容量和费用均不超过 $10^6$，$f \le 10^9$。


\subsection{势能法消除负权边与最短增广路}

最小费用流问题的核心挑战在于：我们不仅需要最大化（或达到指定量的）流量，还需要在诸多合法的流中选择费用最小的那一个。若仅使用基于 BFS 或层次图的增广策略，则完全忽略了边的费用信息，得到的流虽然最大但费用可能远非最优。

解决这一问题的标准框架是 Successive Shortest Path（SSP）算法，其思想朴素而优雅：每一次迭代中，在当前的残量网络中寻找一条从源点到汇点的「最便宜」增广路（即路径上各边费用之和最小的路径），沿该路径输送尽可能多的流量，更新残量网络，重复直至达到需求流量或不存在增广路。这一贪心策略之所以正确，是因为每次选择的都是当前最便宜的可增广路径，而由于反向边的存在，后续的最短路可以「撤销」之前并非最优的输送决策——这一性质与最大流中 Ford--Fulkerson 方法的正确性证明一脉相承。

然而，关键的技术困难在于：残量网络中存在负权边。每当我们在某条边上输送了流量，其反向边的费用为该边原始费用的相反数（负值），用于表示「撤销」流量并获得费用返还。这使得最短路计算无法直接使用 Dijkstra 算法——Dijkstra 要求所有边权非负。一个直观的替代方案是使用 SPFA（Bellman--Ford 的队列优化版），它可以处理负权边，但在每次增广中重复运行 SPFA 的开销（$\mathrm{O}(VE)$ 每次）在流量需求较大时将变得不可接受。

本实现采用了 Johnson 式的势能法（potential method）来消除负权边，从而使 Dijkstra 得以在每一轮迭代中高效运行。核心思想是为每个顶点 $v$ 维护一个势能值 $h[v]$，并将每条边 $(u, v)$ 的原始费用 $w(u, v)$ 替换为约化费用 $w'(u, v) = w(u, v) + h[u] - h[v]$。若势能函数满足三角形不等式 $h[v] \le h[u] + w(u, v)$ 对所有边成立，则所有约化费用非负，Dijkstra 即可安全运行。

初始时，所有原始边费用非负，可设 $h[v] = 0$。每轮 Dijkstra 得到从源点到各顶点的最短约化距离 $d'[v]$ 后，更新势能 $h[v] \leftarrow h[v] + d'[v]$。这一更新保持了约化费用非负的性质（其正确性源于 $d'$ 本身就是最短距离，满足三角不等式），同时使得原始最短路径的实际费用恰好等于更新后的 $h[t]$——这是因为约束 $h[s] = 0$ 始终成立，而沿着最短路径各边的约化费用之和 $d'[t]$ 等于原始费用之和加上 $h[s] - h[t]$，故原始费用即为 $d'[t] + h[t]$。由于势能更新后新的 $h[t]$ 恰好包含了这一轮 Dijkstra 所找到的路径的真实费用，我们可以直接使用 $h[t]$ 作为本轮增广的单位费用。

此后，沿最短路径推送瓶颈流量，更新残量网络，进入下一轮迭代。基于优先队列的 Dijkstra 的单次复杂度为 $\mathrm{O}(E \log V)$，远优于 SPFA 的 $\mathrm{O}(VE)$。

SSP 算法的时间复杂度为 $\mathrm{O}(F \cdot E \log V)$，其中 $F$ 为需求流量（每次增广至少输送 $1$ 单位）。在最坏情况下（例如所有容量均为 $1$ 且须输送大量流量），迭代次数可达 $F$，此时算法效率取决于 $F$ 的大小。然而在本题的数据范围中，$F \le 10^9$ 意味着完全可能出现最坏情况。实践中可以引入容量缩放（capacity scaling）等技术将复杂度降至多项式级别，但对于 $n \le 500$、$m \le 5000$ 的规模，采用势能 Dijkstra 的实现已经具备足够的竞争性。若原始网络中存在负权边（虽然本题样例中均为正），则需要先用一轮 SPFA 或 Bellman--Ford 计算初始势能，以确保约化费用非负。空间复杂度为 $\mathrm{O}(V + E)$。

\subsection{代码实现}

\inputminted{c++}{../src/p4_mincost_flow.cpp}

\section{总结}

本报告以网络流为主题，选取了四道算法思想递进、技术难度层层深入的题目，系统展示了从基础最大流算法到最小费用流扩展的完整知识体系。

题一的 Edmonds--Karp 算法是网络流的入门基石。它以简洁的 BFS 增广替代朴素 Ford--Fulkerson 的随意选路策略，用「最短增广路」这一约束保证了算法在多项式时间内终止。虽然复杂度并非最优，但其思想的清晰性和实现的简便性使其成为理解残量网络、反向边以及增广路这一核心范式的最佳起点。

题二的 Dinic 算法代表了最大流算法的一次质的飞跃。通过「层次图」将一次 BFS 的信息复用为多轮阻塞流推送，Dinic 将 BFS 的调用次数从 $\mathrm{O}(VE)$ 压缩至 $\mathrm{O}(V)$。配合多路增广和当前弧优化，算法在绝大多数实际网络上运行极快——它从「逐条增广」的线性思维跃迁至「分层榨取」的批处理思维，这一思想在后续的预流推进（Push--Relabel）等更高级算法中也有回声。

题三展示了最大流作为一种建模工具的强大威力。二分图最大匹配——一个看似与流毫无关系的组合问题——通过引入超级源汇和单位容量的边，被优雅地归约为最大流问题。更妙的是 K\"onig 定理的构造性运用：残量网络不仅封装了流的信息，其可达性分析还直接给出了最小点覆盖的构造。这提醒我们，残量网络的意义远不止于寻找增广路——它是理解流结构的一扇窗口。

题四将费用维度引入，使问题从「能否输送」拓展为「以最小代价输送」。SSP 算法在概念上延续了增广路的框架，但将 BFS 替换为最短路计算。势能法的引入则体现了算法设计中一项宝贵的技巧：通过维护额外的辅助信息（势能），将一个看似必须忍受负权边的难解情形转化为标准 Dijkstra 可以高效处理的舒适区。这一「重新标定」（re-weighting）的技巧在 Johnson 全源最短路、最小费用流以及许多涉及对偶变量的优化算法中均有深远影响。

回顾四道题目的递进脉络，可以看到两条贯穿始终的方法论线索。其一，残量网络与反向边是流算法的灵魂——它们将「不可逆」的输送操作转化为「可逆」的灵活调配，使得贪心的增广策略在全局上保持最优。其二，算法优化往往源于对冗余计算的消除：Edmonds--Karp 消除随意选路的冗余，Dinic 消除重复 BFS 的冗余，势能法消除负权边带来的 SPFA 冗余——每一次优化都是对问题结构的一次更深刻理解。

\end{document}
